% ===================================================================
%  ਸਾਰਣੀ ਨੰ. 4 — ਭਰੀ ਜਵਾਨੀ ਤੇ ਬੁਢਾਪਾ।
%  Traduction 2026 d'apres texto/io, controlee sur texto/fr, texto/en
%  et texto/hi. Consigne : voir texto/pa/10-sarni-01.tex.
%
%  LA PARENTE SE NOMME PAR LE COTE. Le pendjabi n'a pas un mot pour
%  « beau-pere » mais deux, selon qu'on parle du pere du mari ou de
%  celui de la femme ; de meme pour la belle-sœur — ਨਣਦ est la sœur
%  du mari. On suit l'ido, qui dit toujours de quel cote il parle.
% ===================================================================

%%K t04-apar-1 apar
\VUsaut{4.0mm}
\VUtitre{15.0pt}{60}{ਸਾਰਣੀ ਨੰ. 4}
\VUsaut{1.6mm}
\VUtitre{11.4pt}{0}{ਭਰੀ ਜਵਾਨੀ ਤੇ ਬੁਢਾਪਾ।}
\VUsaut{3.2mm}

%%K t04-tit-1 sub
\VUcentre{11.4pt}{0}{\textit{ਪਹਿਲਾ ਦ੍ਰਿਸ਼।}}
\VUsaut{1.6mm}
\VUtitre{11.4pt}{0}{ਵਿਆਹ ਦੀ ਦਾਵਤ।}
\VUsaut{2.6mm}

%%K t04-tit-2 sub
\VUpk{10.2pt}{40}{ਪੱਤਝੜ।}
\VUsaut{3.0mm}

%%K t04-c1-01-1 p
1. --- ਉਹ ਜਵਾਨ ਹੁਣ ਵੱਡੇ ਹੋ ਗਏ ਹਨ। ਉਨ੍ਹਾਂ ਵਿਚੋਂ ਇਕ, ਜਾਨ, ਵਿਆਹ ਕਰ ਰਿਹਾ ਹੈ। ਅਸੀਂ
\VUgras{ਵਿਆਹ ਦੀ ਦਾਵਤ} ਵਿਚ ਹਾਂ, ਜੋ ਲਾੜੇ ਤੇ ਲਾੜੀ ਦੇ ਦੋਹਾਂ ਪਰਿਵਾਰਾਂ ਨੂੰ ਇਕੋ ਮੇਜ਼
ਦੇ ਚੁਫੇਰੇ ਲੈ ਆਉਂਦੀ ਹੈ।

%%K t04-c1-02-1 p
2. --- ਅਸੀਂ \VUgras{ਪੱਤਝੜ} ਵਿਚ ਹਾਂ, \VUgras{ਸਤੰਬਰ} ਦੇ ਮਹੀਨੇ ਵਿਚ।
\VUgras{ਅਕਤੂਬਰ} ਤੇ \VUgras{ਨਵੰਬਰ} ਦੀ ਠੰਢੀ ਹਵਾ ਨੇ ਹਾਲੇ ਪਿੰਡ ਤੋਂ ਉਸ ਦੇ ਹਰੇ ਪੱਤੇ
ਨਹੀਂ ਖੋਹੇ। ਸ਼ਾਮ ਦੀ ਹਵਾ ਹੌਲੀ ਤੇ ਮਹਿਕਦੀ ਹੈ; ਇਸੇ ਲਈ ਖਾਣਾ ਉਸ \VUgras{ਬਰਾਂਡੇ}
\textsuperscript{(8)} ਦੇ ਹੇਠਾਂ ਵਰਤਾਇਆ ਜਾ ਰਿਹਾ ਹੈ ਜੋ ਬਾਗ਼ ਵੱਲ ਖੁੱਲ੍ਹਦਾ ਹੈ।

%%K t04-c1-03-1 p
3. --- ਦੂਰ \VUgras{ਟਿੱਬੀ} \textsuperscript{(3)} ਉੱਤੇ ਜਾਨ ਦੇ ਮਾਪਿਆਂ ਦਾ ਘਰ ਹੈ, ਇਕ
ਸਜੀਲਾ \VUgras{ਪੇਂਡੂ ਮਕਾਨ} \textsuperscript{(1)}, ਹਰਿਆਵਲ ਵਿਚ ਲੁਕਿਆ ਹੋਇਆ; ਚੰਗੇ
ਅੰਗੂਰਾਂ ਦੇ ਬਾਗ਼, ਜਿਨ੍ਹਾਂ ਤੋਂ ਵਧੀਆ ਸ਼ਰਾਬ ਬਣਦੀ ਹੈ, ਢਲਾਣ ਨੂੰ ਢਕੀ ਬੈਠੇ ਹਨ। ਅੰਗੂਰ
ਉਗਾਉਣ ਵਾਲਾ ਉਨ੍ਹਾਂ ਦੀ ਸੰਭਾਲ ਕਰਦਾ ਹੈ।

%%K t04-c1-04-1 p
4. --- ਵੇਲਾਂ ਦੇ ਉੱਤੋਂ \VUgras{ਤਿੱਤਰਾਂ} \textsuperscript{(2)} ਦਾ ਇਕ ਝੁੰਡ ਉੱਡਦਾ ਹੈ,
ਜੋ \VUgras{ਸ਼ਿਕਾਰ ਦੇ ਰਾਖੇ} \textsuperscript{(5)} ਦੇ ਨੇੜੇ ਆਉਣ ਨਾਲ ਚੌਂਕ ਗਿਆ ਹੈ। ਉਹ,
ਕੱਛ ਵਿਚ \VUgras{ਬੰਦੂਕ} ਦੱਬੀ ਤੇ ਮੋਢੇ ਉੱਤੇ \VUgras{ਸ਼ਿਕਾਰ ਦਾ ਥੈਲਾ} ਲਟਕਾਈ,
\VUgras{ਝਾੜੀਆਂ} \textsuperscript{(4)} ਕੁੱਟ ਰਿਹਾ ਹੈ, ਆਪਣੇ \VUgras{ਕੁੱਤੇ}
\textsuperscript{(6)} ਨਾਲ। ਉਹ ਜਾਗੀਰ ਉੱਤੇ ਨਜ਼ਰ ਰੱਖਦਾ ਹੈ ਤੇ ਚੋਰ-ਸ਼ਿਕਾਰੀਆਂ ਨੂੰ ਸ਼ਿਕਾਰ
ਤਬਾਹ ਨਹੀਂ ਕਰਨ ਦਿੰਦਾ।

%%K t04-c1-05-1 p
5. --- \VUgras{ਬਰਾਂਡੇ} ਦੇ ਬਾਹਰ ਇਕ \VUgras{ਸਿਧਾਈ ਹੋਈ ਵੇਲ} ਚੜ੍ਹਦੀ ਹੈ, ਜਿਸ ਤੋਂ
ਭਾਰੇ \VUgras{ਅੰਗੂਰਾਂ ਦੇ ਗੁੱਛੇ} \textsuperscript{(7)} ਲਟਕਦੇ ਹਨ।

%%K t04-c1-06-1 p
6. --- \VUgras{ਮੇਜ਼} \textsuperscript{(16)} ਨੂੰ ਸਜਾਉਣ ਵਾਲੇ ਪੱਤਝੜ ਦੇ ਸੋਹਣੇ ਫੁੱਲ
\VUgras{ਗੁਲਦਾਊਦੀ} \textsuperscript{(9)} ਹਨ; ਲਾੜੀ ਦੇ \VUgras{ਪਿਤਾ}
\textsuperscript{(10)} ਨੇ ਆਪਣੇ ਕੋਟ ਦੇ ਕਾਜ ਵਿਚ ਇਕ ਲਾਇਆ ਹੋਇਆ ਹੈ।

%%K t04-c1-07-1 p
7. --- ਖਾਣਾ ਮੁੱਕਣ ਉੱਤੇ ਹੈ। \VUgras{ਨੌਕਰ} \textsuperscript{(11)} ਮਿੱਠੇ ਦੇ ਭਾਂਡੇ
ਲਿਆਉਣ ਲੱਗ ਪਿਆ ਹੈ ਤੇ ਥੋੜੀ ਦੇਰ ਪਿੱਛੋਂ ਉਹ ਸਭ ਕੁਝ ਚੁੱਕ ਲੈ ਜਾਵੇਗਾ ਜਿਸ ਦੀ ਹੁਣ ਲੋੜ
ਨਹੀਂ — \VUgras{ਚਮਚੇ} \textsuperscript{(14)}, \VUgras{ਕਾਂਟੇ}
\textsuperscript{(12)}, \VUgras{ਛੁਰੀਆਂ} \textsuperscript{(13)},
\VUgras{ਵਰਤਾਉਣ ਦੇ ਭਾਂਡੇ} \textsuperscript{(15)}।

%%K t04-tit-3 sub
\VUpk{10.2pt}{40}{ਪਰਿਵਾਰ।}
\VUsaut{3.0mm}

%%K t04-c1-08-1 p
8. --- ਸਾਰੇ ਖ਼ੁਸ਼ ਵਿਖਾਈ ਦਿੰਦੇ ਹਨ, ਤੇ ਲਾੜੇ ਦੀ \VUgras{ਮਾਂ} \textsuperscript{(17)}
ਜਿਸ ਮੁਸਕਾਨ ਨਾਲ ਆਪਣੇ ਬੱਚਿਆਂ ਨੂੰ ਵੇਖ ਰਹੀ ਹੈ, ਉਹ ਦੱਸਦੀ ਹੈ ਕਿ ਉਹ ਕਿੰਨੀ ਸੁਖੀ ਹੈ।

%%K t04-c1-09-1 p
9. --- ਇਹ ਵਿਆਹ ਇਲਾਕੇ ਦੇ ਦੋ ਖਾਂਦੇ-ਪੀਂਦੇ ਪਰਿਵਾਰਾਂ ਨੂੰ ਜੋੜਦਾ ਹੈ। ਜਾਨ,
\VUgras{ਪਤੀ} \textsuperscript{(18)}, ਇਕ ਵੱਡੇ ਜ਼ਿਮੀਂਦਾਰ ਦਾ \VUgras{ਪੁੱਤਰ} ਹੈ, ਤੇ
ਉਹ ਇਕ ਕਾਮਯਾਬ ਕਾਰਖ਼ਾਨੇਦਾਰ ਦਾ \VUgras{ਜਵਾਈ} ਬਣ ਰਿਹਾ ਹੈ।

%%K t04-c1-10-1 p
10. --- \VUgras{ਜੋੜੀ} ਜਵਾਨ ਹੈ, ਫਿਰ ਵੀ ਉਨ੍ਹਾਂ ਦੀ \VUgras{ਮੰਗਣੀ} ਬਹੁਤ ਪਹਿਲਾਂ ਹੋ
ਚੁੱਕੀ ਸੀ। \VUgras{ਨਵੀਂ ਲਾੜੀ} \textsuperscript{(19)}, ਮਾਰਥਾ, ਸੰਤਰੇ ਦੇ ਫੁੱਲਾਂ ਨਾਲ
ਸਜੀ ਆਪਣੀ \VUgras{ਚਿੱਟੀ ਪੁਸ਼ਾਕ} ਵਿਚ ਮੋਹ ਲੈਣ ਵਾਲੀ ਲਗਦੀ ਹੈ।

%%K t04-c1-11-1 p
11. --- ਉਸ ਦੇ \VUgras{ਸਹੁਰਾ} \textsuperscript{(20)} ਇੰਨੀ ਚੰਗੀ \VUgras{ਨੂੰਹ} ਲੱਭ
ਕੇ ਖ਼ੁਸ਼ ਹਨ, ਤੇ ਜਾਨ ਦੀ \VUgras{ਸੱਸ} \textsuperscript{(21)} ਆਪਣੀ \VUgras{ਧੀ} ਨੂੰ
ਇੰਨੀ ਸੋਹਣੀ ਵੇਖ ਕੇ ਮੁਸਕਰਾ ਰਹੀ ਹੈ।

%%K t04-c1-12-1 p
12. --- ਮੇਜ਼ ਦੇ ਸਿਰੇ ਉੱਤੇ ਦੂਜੇ \VUgras{ਪ੍ਰਾਹੁਣੇ} \textsuperscript{(22)} ਬੈਠੇ ਹਨ :
\VUgras{ਰਿਸ਼ਤੇਦਾਰ, ਮਿੱਤਰ, ਚਚੇਰੇ ਭਰਾ}। ਇਕ \VUgras{ਮੁੱਖ ਵੇਟਰ} \textsuperscript{(23)},
ਬਾਂਹ ਉੱਤੇ ਰੁਮਾਲ ਪਾਈ, ਫੁਰਤੀ ਨਾਲ ਉਨ੍ਹਾਂ ਦੀ ਸੇਵਾ ਕਰ ਰਿਹਾ ਹੈ।

%%K t04-tit-4 sub
\VUpk{10.2pt}{40}{ਮਿੱਤਰ।}
\VUsaut{3.0mm}

%%K t04-c1-13-1 p
13. --- ਮੇਜ਼ ਦੇ ਸੱਜੇ ਪਾਸੇ ਇਹ ਰਹੀ ਇਕ \VUgras{ਮੁਟਿਆਰ} \textsuperscript{(24)}, ਲਾੜੀ
ਦੀ \VUgras{ਸਹੇਲੀ}, ਫਿਰ ਲਾੜੀ ਦਾ \VUgras{ਭਰਾ}, ਜੋ ਉਸ ਦਾ \VUgras{ਸਰਬਾਲਾ}
\textsuperscript{(25)} ਹੈ। ਉਹ ਆਪਣੀ \VUgras{ਸਰਬਾਲੀ} \textsuperscript{(26)} ਨਾਲ ਗੱਲਾਂ
ਕਰ ਰਿਹਾ ਹੈ, ਜੋ ਲਾੜੇ ਦੀ ਭੈਣ ਹੈ ਤੇ ਇਸ ਵਿਆਹ ਨਾਲ ਮਾਰਥਾ ਦੀ \VUgras{ਨਣਦ} ਬਣ ਰਹੀ ਹੈ।
ਉਹ ਮੋਹ ਲੈਣ ਵਾਲੀ ਮੁਟਿਆਰ ਹੈ। ਉਸ ਕੋਲ ਸੋਹਣੇ \VUgras{ਗਹਿਣੇ} ਹਨ, ਇਕ
\VUgras{ਮੋਤੀਆਂ ਦਾ ਹਾਰ} ਤੇ ਸੋਨੇ ਦਾ ਇਕ \VUgras{ਕੰਗਣ}।

%%K t04-c1-14-1 p
14. --- ਉਨ੍ਹਾਂ ਕੋਲ ਖੜੇ ਹੋ ਕੇ ਗਿਲਾਸ ਚੁੱਕਦੇ ਸੱਜਣ ਪਰਿਵਾਰ ਦੇ ਇਕ \VUgras{ਮਿੱਤਰ}
\textsuperscript{(27)} ਹਨ। ਉਹ \VUgras{ਲਾੜੇ ਦੇ ਗਵਾਹਾਂ} ਵਿਚੋਂ ਇਕ ਹਨ। ਉਹ
\VUgras{ਵਧਾਈ} ਦੇ ਰਹੇ ਹਨ, ਨਵੇਂ ਵਿਆਹਿਆਂ ਦੀ \VUgras{ਸਿਹਤ} ਦੀ ਕਾਮਨਾ ਕਰ ਰਹੇ ਹਨ ਤੇ
ਉਨ੍ਹਾਂ ਨੂੰ ਲੰਮੇ ਸੁਖ ਦੀ ਅਸੀਸ ਦੇ ਰਹੇ ਹਨ। ਉਹ ਪੂਰੇ ਠਾਠ ਵਿਚ ਹਨ : ਲੰਮਾ ਕੋਟ ਤੇ
\VUgras{ਚਮੜੇ ਦੇ ਚਮਕਦੇ ਜੁੱਤੇ} \textsuperscript{(28)}। ਉਹ \VUgras{ਦਾਦੀ}
\textsuperscript{(29)} ਨੂੰ ਨਾਲ ਲਿਆਏ ਹਨ, ਜੋ \VUgras{ਵਿਧਵਾ} ਹੈ।

%%K t04-c1-15-1 p
15. --- \VUgras{ਲੰਮੇ ਕੋਟ} \textsuperscript{(31)} ਵਾਲਾ ਉਹ ਜਵਾਨ, ਜਿਸ ਦੀ ਪਤਲੀ ਕਾਲੀ
ਮੁੱਛ ਹੈ, ਜਾਨ ਦਾ \VUgras{ਭਣਵੱਈਆ} \textsuperscript{(30)} ਹੈ। ਉਹ ਮਸ਼ਹੂਰ ਇੰਜੀਨੀਅਰ ਹੈ।
ਉਹ ਇਕ ਬਹੁਤ ਸਜੀਲੀ \VUgras{ਮੁਟਿਆਰ} \textsuperscript{(33)} ਕੋਲ ਬੈਠਾ ਹੈ, ਜੋ ਮਾਰਥਾ ਦੀ
\VUgras{ਵੱਡੀ ਭੈਣ} ਹੈ ਤੇ ਉਸ ਪਿਆਰੀ ਨਿੱਕੀ ਕੁੜੀ ਦੀ ਮਾਂ, ਜੋ ਆਪਣੇ ਚਚੇਰੇ ਭਰਾ ਦਾ ਹੱਥ
ਫੜੀ ਫੁੱਲ ਦੇਣ ਤੇ \VUgras{« ਸ਼ਾਮ ਨਮਸਕਾਰ »} \VUgras{ਆਖਣ} ਆ ਰਹੀ ਹੈ। ਇਸ ਬੀਬੀ ਦੇ
\VUgras{ਕੰਨ-ਫੁੱਲ} ਬਹੁਤ ਸੋਹਣੇ ਹਨ ਤੇ ਉਸ ਦੇ \VUgras{ਵਾਲਾਂ ਦੇ ਸ਼ਿੰਗਾਰ} ਵਿਚ ਸਲੀਕਾ
ਹੈ।

%%K t04-c1-16-1 p
16. --- ਉਹ ਨਿੱਕਾ ਚਚੇਰਾ ਭਰਾ, ਜੋ ਆਪਣੀ \VUgras{ਕੁੜਤੀ} \textsuperscript{(36)} ਤੇ
ਆਪਣੀ ਢਿੱਲੀ \VUgras{ਨਿੱਕਰ} \textsuperscript{(37)} ਵਿਚ ਇੰਨਾ ਅਨੋਖਾ ਲਗਦਾ ਹੈ, ਬੜਾ ਤਰਸ
ਜੋਗ ਹੈ। ਉਹ \VUgras{ਯਤੀਮ} \textsuperscript{(35)} ਹੈ। ਬਹੁਤ ਨਿੱਕੀ ਉਮਰ ਵਿਚ ਉਸ ਨੇ ਆਪਣਾ
ਪਿਤਾ ਤੇ ਆਪਣੀ ਮਾਂ ਗੁਆ ਲਏ। ਉਸ ਦੇ ਚਾਚਾ ਉਸ ਦੇ \VUgras{ਸਰਪ੍ਰਸਤ} \textsuperscript{(38)}
ਹਨ, ਤੇ ਉਸ ਵਿਚਾਰੇ ਬਾਲ ਨੂੰ ਪਾਲਣ ਲਈ ਹੀ ਛੜੇ ਰਹਿ ਗਏ। ਉਹ ਭਲੇ ਬੰਦੇ ਹਨ। ਉਨ੍ਹਾਂ ਦੇ ਸਿਰ
ਉੱਤੇ ਵਾਲ ਘੱਟ ਹਨ ਤੇ ਉਨ੍ਹਾਂ ਨੂੰ ਦੂਰ ਦਾ ਬਹੁਤ ਘੱਟ ਦਿਸਦਾ ਹੈ; ਉਹ \VUgras{ਨੱਕ-ਐਨਕ}
ਲਾਉਂਦੇ ਹਨ।

%%K t04-tit-5 sub
\VUsaut{2.4mm}
\VUpk{10.2pt}{40}{ਮਿੱਠਾ।}
\VUsaut{3.0mm}

%%K t04-c1-17-1 p
17. --- ਪ੍ਰਾਹੁਣਿਆਂ ਦੇ ਪਿੱਛੇ, \VUgras{ਮੇਜ਼ਪੋਸ਼} ਵਿਛੀ ਇਕ ਮੇਜ਼ ਉੱਤੇ, \VUgras{ਮਿੱਠਾ}
\textsuperscript{(39)} ਸਜਿਆ ਹੋਇਆ ਹੈ। ਇਕ ਵੱਡੇ ਕਟੋਰੇ ਵਿਚ ਇਹ ਰਹੇ ਫਲ :
\VUgras{ਆੜੂ} \textsuperscript{(40)}, \VUgras{ਨਾਸ਼ਪਾਤੀ} \textsuperscript{(41)},
\VUgras{ਸੇਬ} \textsuperscript{(42)}, \VUgras{ਖ਼ੁਬਾਨੀ} \textsuperscript{(43)} ਤੇ
\VUgras{ਅੰਗੂਰ} \textsuperscript{(44)}। ਉਨ੍ਹਾਂ ਕੋਲ \VUgras{ਅਨਾਨਾਸ}
\textsuperscript{(45)}, ਇਕ ਸੋਹਣਾ \VUgras{ਕੇਕ} \textsuperscript{(46)},
\VUgras{ਸ਼ਰਾਬ ਦੀਆਂ ਬੋਤਲਾਂ} \textsuperscript{(48)} ਤੇ \VUgras{ਸ਼ੈਂਪੇਨ}
\textsuperscript{(49)}, ਤੇ ਸਾਫ਼ \VUgras{ਪਲੇਟਾਂ} \textsuperscript{(47)} ਦੇ ਢੇਰ ਵੀ ਹਨ।

%%K t04-c1-18-1 p
18. --- \VUgras{ਸ਼ਰਾਬ ਵਰਤਾਉਣ ਵਾਲਾ} \textsuperscript{(50)} ਪੁਰਾਣੀ ਸ਼ਰਾਬ ਦੀ ਇਕ
\VUgras{ਬੋਤਲ} \textsuperscript{(52)} ਖੋਲ੍ਹ ਰਿਹਾ ਹੈ, \VUgras{ਪੇਚ}
\textsuperscript{(51)} ਨਾਲ। \VUgras{ਡਾਟ} \textsuperscript{(53)} ਬੜੀ ਔਖ ਨਾਲ ਨਿਕਲਦੀ
ਲਗਦੀ ਹੈ। ਇਸ ਵਰਤਾਉਣ ਵਾਲੇ ਦੀ ਬਾਂਹ ਉੱਤੇ ਇਕ \VUgras{ਰੁਮਾਲ} \textsuperscript{(54)} ਹੈ।

%%K t04-c1-19-1 p
19. --- ਖਾਣੇ ਪਿੱਛੋਂ ਸੱਜਣ ਲੋਕ ਸਿਗਰਟ ਵਾਲੇ ਕਮਰੇ ਵਿਚ ਜਾਣਗੇ, ਤੇ ਬੀਬੀਆਂ ਨਾਚ ਲਈ
ਤਿਆਰ ਹੋਣ ਤੁਰ ਪੈਣਗੀਆਂ।

%%K t04-tit-6 sub
\VUsaut{5.0mm}
\VUcentre{11.4pt}{0}{\textit{ਦੂਜਾ ਦ੍ਰਿਸ਼।}}
\VUsaut{1.6mm}
\VUpk{11.4pt}{40}{ਦਾਦੇ ਦਾ ਜਨਮ-ਦਿਨ।}
\VUsaut{2.6mm}

%%K t04-tit-7 sub
\VUpk{10.2pt}{40}{ਸੌਗਾਤ।}
\VUsaut{3.0mm}

%%K t04-c2-01-1 p
1. --- ਦਸ ਵਰ੍ਹੇ ਲੰਘ ਗਏ। ਜਾਨ ਦੇ ਮਾਪੇ ਬੁੱਢੇ ਹੋ ਚੱਲੇ ਹਨ, ਤੇ ਉਹ ਆਪਣੇ ਤਿੰਨਾਂ
ਪੋਤਰਿਆਂ ਨੂੰ ਬਹੁਤ ਪਿਆਰ ਕਰਦੇ ਹਨ — ਦੋ \VUgras{ਪੋਤਰੇ} \textsuperscript{(2)} ਤੇ ਇਕ
\VUgras{ਪੋਤਰੀ} \textsuperscript{(3)}। ਅੱਜ \VUgras{ਦਾਦੇ ਦਾ ਜਨਮ-ਦਿਨ} ਹੈ, ਇਸੇ ਲਈ
ਬੱਚੇ ਉਨ੍ਹਾਂ ਨੂੰ ਵਧਾਈ ਦੇਣ ਆਏ ਹਨ।

%%K t04-c2-02-1 p
2. --- ਨਿੱਕਾ ਪੀਟਰ ਹਾਲੇ ਬਹੁਤ ਨਿੱਕਾ ਹੈ; ਉਹ ਸਿਰਫ਼ ਤਿੰਨ ਵਰ੍ਹਿਆਂ ਦਾ ਹੈ। ਉਹ
\VUgras{ਬਿੱਲੀ} \textsuperscript{(1)} ਨੂੰ ਆਪਣੇ ਵੱਲ ਆਉਂਦੀ ਵੇਖਦਾ ਹੈ। ਉਹ ਉਸ ਦੀ ਸੋਹਣੀ
ਰੇਸ਼ਮੀ \VUgras{ਜੱਤ} ਤੇ ਲੰਮੀ \VUgras{ਪੂਛ} ਉੱਤੇ ਹੱਥ ਫੇਰਨਾ ਚਾਹੁੰਦਾ ਹੈ; ਉਸ ਨੂੰ ਇਹ
ਨਹੀਂ ਸੁੱਝਦਾ ਕਿ ਉਨ੍ਹਾਂ ਸੋਹਣੇ \VUgras{ਪੰਜਿਆਂ} ਦੇ ਹੇਠਾਂ ਤਿੱਖੇ ਨਹੁੰ ਹਨ, ਜੋ ਉਸ ਨੂੰ
ਲਿਤਾੜ ਵੀ ਸਕਦੇ ਹਨ।

%%K t04-c2-03-1 p
3. --- ਉਸ ਦੀ ਭੈਣ ਗੈਬਰੀਏਲ, ਜੋ ਉਸ ਦਾ ਹੱਥ ਫੜੀ ਹੋਈ ਹੈ, ਕਿਤੇ ਵੱਧ ਗੰਭੀਰ ਹੈ, ਤੇ
ਮਾਰਸਲ, ਆਪਣੇ ਵੱਡੇ \VUgras{ਓਵਰਕੋਟ} \textsuperscript{(4)} ਤੇ ਚਮੜੇ ਦੀ \VUgras{ਪੇਟੀ}
\textsuperscript{(5)} ਵਿਚ, ਪੂਰਾ ਨਿੱਕਾ ਬੰਦਾ ਲਗਦਾ ਹੈ, ਭਾਵੇਂ ਨਿੱਕੀ ਨਿੱਕਰ ਉਸ ਦੀਆਂ
\VUgras{ਜੁਰਾਬਾਂ} \textsuperscript{(6)} ਖੁੱਲ੍ਹੀਆਂ ਛੱਡ ਦਿੰਦੀ ਹੈ।

%%K t04-c2-04-1 p
4. --- ਉਨ੍ਹਾਂ ਦੇ ਪਿਤਾ ਨੇ ਆਪਣਾ ਮੋਟਾ ਸਿਆਲ ਦਾ \VUgras{ਓਵਰਕੋਟ}
\textsuperscript{(7)} ਪਾਇਆ ਹੋਇਆ ਹੈ, \VUgras{ਜੱਤ ਵਾਲੇ ਕਾਲਰ} \textsuperscript{(8)}
ਸਮੇਤ, ਤੇ ਆਪਣੀ \VUgras{ਜੇਬ} \textsuperscript{(9)} ਵਿਚ ਇਕ ਗਲੂਬੰਦ ਰੱਖਿਆ ਹੋਇਆ ਹੈ।
ਉਨ੍ਹਾਂ ਦੀ ਘਰ ਵਾਲੀ ਨੇ \VUgras{ਖੰਭਾਂ} \textsuperscript{(10)} ਨਾਲ ਸਜੀ ਨਮਦੇ ਦੀ ਹੈਟ,
ਆਪਣੀ \VUgras{ਜੈਕਟ} \textsuperscript{(11)}, ਆਪਣਾ ਜੱਤ ਵਾਲਾ \VUgras{ਗਲੂਬੰਦ}
\textsuperscript{(12)} ਤੇ ਆਪਣਾ \VUgras{ਮਫ਼} \textsuperscript{(13)} ਪਾਇਆ ਹੋਇਆ ਹੈ।

%%K t04-c2-05-1 p
5. --- ਬਹੁਤ ਪਾਲਾ ਹੈ, ਕਿਉਂਕਿ ਸਿਆਲ ਆ ਗਿਆ ਹੈ। \VUgras{ਬਰਫ਼} \textsuperscript{(19)}
ਸਾਰਾ ਦਿਨ ਪੈਂਦੀ ਰਹੀ ਹੈ ਤੇ ਘਰਾਂ ਦੀਆਂ ਛੱਤਾਂ ਬਿਲਕੁਲ ਚਿੱਟੀਆਂ ਹਨ। \VUgras{ਰਾਤ} ਨਾਲ
ਪਾਲਾ ਹੋਰ ਵੀ ਤਿੱਖਾ ਹੋ ਜਾਂਦਾ ਹੈ। ਅਸਮਾਨ ਵਿਚ \VUgras{ਚੰਨ} \textsuperscript{(17)} ਤੇ
\VUgras{ਤਾਰੇ} \textsuperscript{(18)} ਚਮਕਦੇ ਹਨ ਤੇ \VUgras{ਹਨੇਰੇ} ਨੂੰ ਚੀਰਦੇ ਹਨ।

%%K t04-tit-8 sub
\VUsaut{4.2mm}
\VUpk{10.2pt}{40}{ਬੀਮਾਰੀ।}
\VUsaut{3.0mm}

%%K t04-c2-06-1 p
6. --- ਉਹ ਵਿਚਾਰਾ \VUgras{ਅੰਨ੍ਹਾ} \textsuperscript{(20)}, ਜੋ \VUgras{ਦਵਾਖ਼ਾਨੇ}
\textsuperscript{(16)} ਦੇ ਸਾਹਮਣੇ ਚੌਕ ਲੰਘ ਰਿਹਾ ਹੈ ਤੇ ਜਿਸ ਨੂੰ ਉਸ ਦਾ
\VUgras{ਪੂਡਲ} \textsuperscript{(21)} ਰਾਹ ਵਿਖਾ ਰਿਹਾ ਹੈ, ਬੜਾ ਨਿਘਰਿਆ ਹੈ।
\VUgras{ਨੌਕਰਾਣੀ} \textsuperscript{(14)} ਜਸਟੀਨ ਰਸੋਈ ਤੋਂ ਬਹੁਤ ਗਰਮ \VUgras{ਕਾੜ੍ਹੇ}
ਦਾ ਇਕ \VUgras{ਕਟੋਰਾ} \textsuperscript{(15)} ਲਿਆ ਰਹੀ ਹੈ, ਖ਼ੂਬ ਮਿੱਠਾ ਕੀਤਾ ਹੋਇਆ।

%%K t04-c2-07-1 p
7. --- \VUgras{ਦਾਦੀ} \textsuperscript{(23)}, ਜੋ ਮੇਜ਼ ਉੱਤੋਂ ਆਪਣੀ \VUgras{ਐਨਕ}
\textsuperscript{(26)} ਚੁੱਕਣਾ ਚਾਹੁੰਦੀ ਹੈ ਤਾਂ ਜੋ ਉਹ \VUgras{ਪਰਚੀ}
\textsuperscript{(27)} ਪੜ੍ਹ ਸਕੇ ਜੋ \VUgras{ਵੈਦ} \textsuperscript{(25)} ਨੇ ਹੁਣੇ ਲਿਖੀ
ਹੈ, ਆਪਣੀ \VUgras{ਸੋਟੀ} \textsuperscript{(24)} ਦੇ ਸਹਾਰੇ ਆਪਣੀ ਕੁਰਸੀ ਤੋਂ ਉੱਠ ਰਹੀ ਹੈ।

%%K t04-c2-08-1 p
8. --- ਉਨ੍ਹਾਂ ਨੂੰ ਅਕਸਰ ਨਿੱਕੀਆਂ ਨਿੱਕੀਆਂ \VUgras{ਤਕਲੀਫ਼ਾਂ} ਰਹਿੰਦੀਆਂ ਹਨ,
\VUgras{ਚੱਕਰ}, \VUgras{ਜ਼ੁਕਾਮ} ਤੇ \VUgras{ਦੰਦ ਦੀ ਪੀੜ}; ਉਨ੍ਹਾਂ ਦੀਆਂ ਲੱਤਾਂ ਕਮਜ਼ੋਰ
ਹਨ ਤੇ ਉਨ੍ਹਾਂ ਦੀਆਂ ਅੱਖਾਂ ਹੁਣ ਚੰਗੀਆਂ ਨਹੀਂ ਰਹੀਆਂ। ਫਿਰ ਵੀ ਉਹ ਆਪਣੇ ਬੁੱਢੇ
\VUgras{ਵੈਦ} ਨੂੰ ਛੇੜ ਕੇ ਖ਼ੁਸ਼ ਹੁੰਦੀ ਹੈ, ਜੋ ਉਨ੍ਹਾਂ ਨੂੰ ਸਦਾ ਕੋਈ \VUgras{ਘੁੱਟ}
\textsuperscript{(28)} ਲਿਖ ਦੇਣਾ ਚਾਹੁੰਦੇ ਹਨ। ਅੱਜ ਉਹ ਆਪਣਾ ਜਵਾਨ ਪਰਿਵਾਰ ਆਪਣੇ ਚੁਫੇਰੇ
ਲੱਭ ਕੇ ਬਹੁਤ ਸੁਖੀ ਹੈ।

%%K t04-c2-09-1 p
9. --- ਚੰਗੀ ਤਰ੍ਹਾਂ ਬੰਦ ਕਮਰੇ ਵਿਚ ਨਿੱਘ ਹੈ। ਸੰਗਮਰਮਰ ਦੀ \VUgras{ਅੰਗੀਠੀ}
\textsuperscript{(29)} ਵਿਚ \VUgras{ਸੋਹਣੀ ਅੱਗ} \textsuperscript{(30)} ਬਲ ਰਹੀ ਹੈ, ਤੇ
ਹੁਣੇ ਹੁਣੇ ਬਾਲੇ \VUgras{ਲੈਂਪ} \textsuperscript{(31)} ਦੀ ਕੋਮਲ \VUgras{ਲੋਅ} ਵਿਚ ਦਿਲ
ਖ਼ੁਸ਼ ਹੋ ਜਾਂਦਾ ਹੈ।

%%K t04-tit-9 sub
\VUsaut{4.2mm}
\VUpk{10.2pt}{40}{ਦਾਦਾ।}
\VUsaut{3.0mm}

%%K t04-c2-10-1 p
10. --- \VUgras{ਦਾਦਾ} \textsuperscript{(33)} ਤਕ ਭੁੱਲ ਗਏ ਹਨ ਕਿ ਉਹ ਬੀਮਾਰ ਰਹੇ ਹਨ, ਕਿ
ਉਨ੍ਹਾਂ ਦਾ \VUgras{ਜੋੜਾਂ ਦਾ ਦਰਦ} ਉਨ੍ਹਾਂ ਨੂੰ ਆਪਣੀ \VUgras{ਆਰਾਮ-ਕੁਰਸੀ}
\textsuperscript{(34)} ਵਿਚ ਬੰਨ੍ਹੀ ਰੱਖਦਾ ਹੈ, ਤੇ ਕਿ ਕੱਲ੍ਹ ਹੀ ਉਨ੍ਹਾਂ ਨੂੰ
\VUgras{ਤਾਪ ਤੇ ਗ਼ਸ਼} ਆਈ ਸੀ। ਉਨ੍ਹਾਂ ਨੂੰ ਹੁਣ ਯਾਦ ਨਹੀਂ ਕਿ ਪਿਛਲੇ ਸਿਆਲ ਉਨ੍ਹਾਂ ਨੂੰ
\VUgras{ਨਮੂਨੀਆ} ਹੋਇਆ ਸੀ, ਤੇ ਇਕ ਸਖ਼ਤ ਪੀੜ ਦੇਣ ਵਾਲੀ \VUgras{ਖੰਘ} ਨੇ ਉਨ੍ਹਾਂ ਨੂੰ ਬਹੁਤ
ਸਤਾਇਆ ਸੀ।

%%K t04-c2-10-2 p
ਤੇ ਉਹ ਔਖ ਨਾਲ ਹੀ ਰਾਜ਼ੀ ਹੋਏ। ਹੁਣ ਕੁਝ ਚੰਗੇ ਹਨ। ਪਰ ਉਨ੍ਹਾਂ ਦੀ ਲੱਤ, ਜਿਸ ਨੂੰ ਉਹ ਇਕ
\VUgras{ਗੱਦੀ} \textsuperscript{(38)} ਉੱਤੇ ਟਿਕਾਈ ਰੱਖਦੇ ਹਨ, ਤੇ ਉਹ ਗੱਦੀ ਇਕ ਨਿੱਕੀ
\VUgras{ਚੌਕੀ} \textsuperscript{(39)} ਉੱਤੇ ਰੱਖੀ ਹੈ, ਹੁਣ ਵੀ ਦੁਖਦੀ ਹੈ।

%%K t04-c2-11-1 p
11. --- ਜਦੋਂ ਉਹ ਆਪਣੇ ਨਿੱਘੇ \VUgras{ਡਰੈਸਿੰਗ ਗਾਊਨ} \textsuperscript{(35)} ਵਿਚ, ਆਪਣੇ
ਵੱਡੇ ਸਿੱਪੀ ਦੇ \VUgras{ਬਟਨਾਂ} \textsuperscript{(36)} ਸਮੇਤ, ਚੰਗੀ ਤਰ੍ਹਾਂ ਲਪੇਟੇ ਹੁੰਦੇ
ਹਨ, ਤੇ ਜਦੋਂ ਉਨ੍ਹਾਂ ਕੋਲ ਅਖ਼ਬਾਰ ਪੜ੍ਹਨ ਲਈ ਉਹ ਨਿੱਕੀ \VUgras{ਯਤੀਮ ਕੁੜੀ}
\textsuperscript{(40)} ਹੁੰਦੀ ਹੈ, ਆਪਣੀ ਚਿੱਟੀ \VUgras{ਟੋਪੀ}, ਆਪਣੀ ਨੀਲੀ
\VUgras{ਪੁਸ਼ਾਕ} \textsuperscript{(42)} ਤੇ \VUgras{ਚਾਦਰ} \textsuperscript{(41)} ਵਿਚ,
ਤਦੋਂ ਉਹ ਕੋਈ ਸ਼ਿਕਾਇਤ ਨਹੀਂ ਕਰਦੇ।

%%K t04-c2-12-1 p
12. --- ਹਰ ਵਰ੍ਹੇ, ਜਦੋਂ \VUgras{ਜੰਤਰੀ} \textsuperscript{(43)} ਪਹਿਲੀ
\VUgras{ਦਸੰਬਰ} ਦੀ \VUgras{ਤਾਰੀਖ਼} ਮੁੜ ਲੈ ਆਉਂਦੀ ਹੈ, ਉਹ ਆਪਣੇ \VUgras{ਪੋਤਰਿਆਂ} ਦੀ
ਬੇਸਬਰੀ ਨਾਲ ਉਡੀਕ ਕਰਦੇ ਹਨ, ਕਿਉਂਕਿ ਉਹ ਜਾਣਦੇ ਹਨ ਕਿ ਉਹ ਉਸ ਅਹਿਮ \VUgras{ਤਾਰੀਖ਼} ਨੂੰ
ਨਹੀਂ ਭੁੱਲਣਗੇ।

%%K t04-c2-13-1 p
13. --- ਉਹ ਭਾਵੁਕ ਹੋ ਕੇ ਉਹ ਦੂਰ ਦੇ ਦਿਨ ਯਾਦ ਕਰਦੇ ਹਨ ਜਦੋਂ ਉਹ ਆਪ ਉਸ ਸ਼ਾਨਦਾਰ ਸ਼ਾਹੀ
\VUgras{ਮਾਰਸ਼ਲ} ਨੂੰ ਵਧਾਈ ਦੇਣ ਜਾਇਆ ਕਰਦੇ ਸਨ, ਜਿਨ੍ਹਾਂ ਦੀ \VUgras{ਤਸਵੀਰ}
\textsuperscript{(44)} ਕੰਧ ਉੱਤੇ ਟੰਗੀ ਹੈ, ਤੇ ਜੋ ਬੱਚਿਆਂ ਦੇ \VUgras{ਪੜਦਾਦਾ} ਹਨ।
\VUsaut{6.0mm}
\VUfilet{22mm}
